\chapter{Conclusion and Perspectives}
\label{chap:conclusion}

\section{Project Assessment}

This internship set out to build a working, secure civic reporting platform for Tunisian municipalities, and delivered exactly that: a trilingual mobile application, a municipal staff dashboard, and a backend API, all communicating over a REST interface backed by a geospatial relational database. Beyond the functional core, two additional outcomes turned out to be at least as significant as the original scope: a full, severity-ranked security audit that found and fixed a genuinely exploitable set of issues (a retrievable signing key, unauthenticated sensitive endpoints, non-revocable sessions), and a continuous integration/deployment pipeline that itself surfaced real, previously-undetected defects --- a mobile smoke test that had silently never passed, and several dependencies with known vulnerabilities that had drifted unnoticed.

If there is one overall lesson from the internship, it is that \textbf{verification changes outcomes, not just confidence}: more than one fix in Chapter~\ref{chap:realisation} was incomplete on its first attempt and was only caught by a subsequent, independent check. Treating audit, fix, and verification as separate, repeated steps --- rather than folding verification into ``I fixed it, moving on'' --- was the single methodological choice that mattered most.

\section{Personal and Technical Contributions}

On the technical side, this internship required working across the full stack of a real product: a Flutter mobile application with offline support and trilingual RTL-aware UI, a React/TypeScript dashboard, a FastAPI backend with a PostGIS-backed geospatial data model, and a GitHub Actions CI/CD pipeline with real service-container-based integration tests. It also required a security-specific skill set distinct from feature development: reasoning about attacker-visible attack surface (what can an unauthenticated request reach?), about race conditions in concurrent request handling, and about the operational side of key rotation and credential hygiene in a deployed system, not just a textbook description of them.

On the personal side, working with an LLM-assisted development workflow (Section~\ref{chap:etude}) required developing a specific discipline: treating the agent's output as a proposal to verify, not a finished answer to trust, and building that verification into the workflow itself (a curl-based regression check, an automated test, an independent review pass) rather than relying on visual code inspection alone.

\section{Possible Improvements and Perspectives}

The following items were identified during the internship as genuine next steps, not filler:

\begin{itemize}[leftmargin=*]
  \item \textbf{Municipality-scoped staff visibility} --- reports now auto-populate their nearest municipality, but non-admin staff roles still see every report platform-wide rather than only their own municipality's. This is the most immediately impactful next feature.
  \item \textbf{AI classification microservice} --- the data model already reserves fields for AI-suggested category, confidence score, and duplicate detection; the model itself was out of scope for this internship and remains to be built.
  \item \textbf{Load testing} --- formal concurrent-load testing against the deployed environment, to validate the sub-2-second response target under realistic traffic rather than only during manual single-user testing.
  \item \textbf{Git history remediation} --- purging the historical commit containing the now-rotated signing key from the repository's history entirely, deferred as a lower-urgency hygiene task once the key itself was rotated and the repository was made private.
  \item \textbf{Upload content validation} --- the current upload path validates a file's declared content type; validating the actual file signature (``magic bytes'') server-side would close a narrow gap where a client could mislabel an upload's type.
  \item \textbf{Email verification persistence} --- the email-verification endpoint currently confirms a one-time code but does not yet persist a verified/unverified state on the account, since nothing in the current product gates behavior on it; this is a product decision to make before it becomes a schema change.
\end{itemize}

Taken together, these are consistent with a platform that reached a genuinely working, secure, and continuously-tested state within the internship's timeframe, with a clear and specific --- rather than vague --- list of what comes next.
